\section{Estado del arte}
\label{sec:estado_arte}

Los trabajos más cercanos a esta tesina se agrupan en tres ejes: la clasificación de lesiones de piel, los
\glspl{vlm} para dermatología y la generación de reportes dermatológicos. La revisión se concentra en trabajos
publicados desde 2020, y los anteriores aparecen solo como hitos, con su fecha. Varios de 2025 y 2026 están
publicados solo en arXiv o medRxiv, sin revisión de pares, y en cada caso se lo indica.

\subsection{Diagnóstico automático de lesiones de piel}
\label{sec:estado_arte_diagnostico}

La clasificación de lesiones asigna a una imagen una entre varias categorías diagnósticas, sobre imágenes
dermatoscópicas (HAM10000 y los desafíos ISIC, conjuntos de referencia muy utilizados) o fotografías clínicas.
Importa cómo trata cada trabajo el desbalance de clases, con qué métrica elige el modelo, si mantiene las imágenes
de una misma lesión en una partición, si valida fuera de la fuente y si compara contra clínicos. La
Tabla~\ref{tab:marco_sota_clasificacion} resume los trabajos.

\begin{table}[!t]
    \centering
    \caption{Trabajos de clasificación de lesiones de piel. PCP: médicos de atención primaria; NP: enfermeros
    especializados (\emph{nurse practitioners}).}
    \label{tab:marco_sota_clasificacion}
    \footnotesize
    \setlength{\tabcolsep}{3pt}
    \begin{tabularx}{\textwidth}{@{}>{\raggedright\arraybackslash}p{1.9cm}
        c
        >{\raggedright\arraybackslash\hspace{0pt}}X >{\raggedright\arraybackslash\hspace{0pt}}X
        >{\raggedright\arraybackslash\hspace{0pt}}p{1.8cm} >{\raggedright\arraybackslash\hspace{0pt}}X
        >{\raggedright\arraybackslash\hspace{0pt}}p{1.75cm}@{}}
        \toprule
        \textbf{Trabajo} & \textbf{Año} & \textbf{Arquitectura} & \textbf{Datos} & \textbf{Métricas} & \textbf{Resultado principal} & \textbf{Expertos} \\
        \midrule
        \textcite{liu2020deep} & 2020 & Ensamble de 6 Inception-v4 con metadatos
            & 17\,777 casos de 17 sitios; 26 condiciones
            & Exactitud, sensibilidad promedio
            & Top-1 0.67; dermatólogos 0.63, PCP 0.45, NP 0.41
            & 18 lectores; 3 dermatólogos de referencia \\
        \filasep
        \textcite{gessert2020skin} & 2020 & Ensamble de CNN con metadatos
            & ISIC 2019; 5 particiones por lesión
            & Sensibilidad macro, AUC
            & Sin metadatos: 72.5\,\% en validación cruzada, 63.6\,\% en prueba
            & No \\
        \filasep
        \textcite{yao2022single} & 2022 & RegNetY única con pérdida propia
            & ISIC 2017, 2018, 2019 y 7-point
            & Sensibilidad macro y de melanoma, AUC
            & ISIC 2018: 0.875 y 0.819; ISIC 2019: 0.645
            & No \\
        \filasep
        \textcite{marchetti2023prospective} & 2023 & Ensamble de 18 modelos
            & 603 lesiones biopsiadas, 95 melanomas
            & Sensibilidad, especificidad, AUC
            & Sensibilidad 96.8\,\%; especificidad 37.4\,\%; AUC 0.857
            & 11 dermatólogos; impacto teórico \\
        \filasep
        \textcite{kreouzi2024learning} & 2024 & Cuatro CNN preentrenadas, entre ellas DenseNet121
            & DermNet; melanoma contra nevo
            & Exactitud, AUC, tiempo de inferencia, tamaño
            & Exactitud 92.30\,\% (DenseNet121); AUC 0.957 (ResNet50V2)
            & No \\
        \filasep
        \textcite{aboulmira2024skinhealthmate} & 2024 & Ensamble de EfficientNetB1 y B5 en una aplicación
            & HAM10000; 7 clases
            & Exactitud, F1 macro, AUC
            & Exactitud 93\,\%; F1 macro 0.81; sensibilidad de melanoma 0.49
            & No \\
        \filasep
        \textcite{khan2025ensemble} & 2025 & Ensamble de tres CNN, más LLM
            & Nevo contra carcinoma basocelular; 80/10/10 al azar
            & Exactitud, F1, pérdida logarítmica
            & Exactitud 97.67\,\% contra 93.83\,\% (ResNet50)
            & No \\
        \filasep
        \textcite{kim2025skinsavvy2} & 2025 & Swin o EfficientNetV2, más GPT
            & HAM10000 ampliado con imágenes externas
            & Exactitud, precisión, sensibilidad, F1
            & Exactitud balanceada 0.780; F1 promedio 0.793 (Swin)
            & 2 dermatólogos entre 20 evaluadores \\
        \bottomrule
    \end{tabularx}
\end{table}

Dos hitos fijaron la tarea como un terreno en el que el aprendizaje profundo compite con el especialista. En 2017
una \gls{cnn} igualó a dermatólogos certificados en dos distinciones entre lesiones malignas y benignas, carcinoma
queratinocítico contra queratosis seborreica y melanoma contra nevo~\cite{esteva2017dermatologistlevel}. En 2019,
los clasificadores del desafío ISIC 2018, entrenados con HAM10000, acertaron en su prueba en promedio más que 511 participantes humanos, y los tres mejores superaron también a los expertos con
más de diez años de experiencia~\cite{tschandl2019comparison}.

\textcite{liu2020deep} extendieron la tarea de la lesión aislada a las 26 condiciones de piel más frecuentes en
atención primaria, a partir de fotografías clínicas y de la historia del paciente. Un ensamble de seis redes
Inception-v4 fusiona las imágenes de cada caso con 45 variables de metadatos y se entrena con etiquetas blandas
construidas con los votos de un panel de dermatólogos. Los datos son casos de teledermatología de 17 sitios de
Estados Unidos, con partición temporal y separada por paciente, y la referencia es el diagnóstico de tres dermatólogos, sin confirmación
patológica. Se mide la exactitud y la sensibilidad promedio entre condiciones. Frente a 18 clínicos, en un
subconjunto enriquecido con condiciones raras, alcanzó una exactitud top-1 de 0.67 contra 0.63 de los dermatólogos y 0.45
de los médicos de atención primaria; entre los tumores malignos, en cambio, su sensibilidad top-1 fue 0.57 contra
0.71 (0.88 contra 0.89 entre los tres primeros diagnósticos). Los
autores reconocen que la partición temporal no reemplaza a un conjunto externo.

\textcite{gessert2020skin} describen el ensamble de EfficientNet, SENet154 y ResNeXt que ganó las dos tareas del
desafío ISIC 2019. Ponderan la entropía cruzada por la frecuencia inversa de cada clase y agregan una novena clase,
con imágenes externas, para la clase desconocida de la prueba. Parten de 25\,331 imágenes, más datos externos, en cinco
particiones que mantienen juntas las de una misma lesión. La métrica es la sensibilidad macro, y la evidencia es el primer puesto, sin
comparación con otros métodos. En la tarea sin metadatos llegan a 72.5\,\% en validación cruzada, pero a 63.6\,\% en la
prueba oficial, una brecha que atribuyen en parte a la clase desconocida; con metadatos, 74.2\,\% y 63.4\,\%. No hubo expertos.

\textcite{yao2022single} proponen en cambio un modelo único, de bajo costo, para conjuntos pequeños y
desbalanceados. De 41 \glspl{cnn} eligen una RegNetY y le suman una pérdida con un coeficiente por clase, mayor para
el melanoma, y un aprendizaje acumulativo que pasa del entrenamiento convencional al rebalanceo. Usan ISIC 2017, 2018,
2019 y 7-point, con la sensibilidad macro como métrica principal y la del melanoma aparte, y se comparan con los
ganadores de los desafíos. En ISIC 2019 alcanzan una sensibilidad macro de 0.645 contra 0.636 del ganador original del
desafío, y en ISIC 2018 muestran que subir el coeficiente del melanoma eleva su sensibilidad de 0.772 a 0.871 a
costa de la sensibilidad macro. No
hubo expertos. Reconocen que los ensambles rinden en general mejor que un modelo único. El artículo no describe una
partición de validación: la selección se hace sobre la prueba de ISIC 2018.

\textcite{marchetti2023prospective} llevaron fuera del banco de pruebas el ganador del desafío SIIM-ISIC 2020,
un ensamble de 18 modelos EfficientNet y ResNet; cada modelo se entrenó en cinco
particiones, de modo que para cada imagen se ejecutan y promedian 90 redes. Lo validaron de forma
prospectiva, con biopsia como referencia, en un estudio observacional con 603 lesiones, 95 de ellas melanomas. Con
el umbral fijado de antemano, la sensibilidad fue 96.8\,\% y la especificidad 37.4\,\%, y el \gls{auc} de 0.857 fue
similar al de la prueba del desafío y superior al de los dermatólogos sin ayuda (0.780). Con acceso al puntaje, el
\gls{auc} de 11 dermatólogos subió a 0.816, aunque el impacto es teórico porque no se modificó la atención de ningún
paciente. Los autores reconocen que fue un solo centro, con una muestra pequeña y sin lesiones biopsiadas de pacientes de
fototipo V o VI.

Los trabajos recientes sobre HAM10000 y conjuntos afines se comparan casi siempre entre arquitecturas, sin expertos
ni un protocolo común. \textcite{kreouzi2024learning} enfrentan cuatro \glspl{cnn} (DenseNet121, ResNet50V2,
NASNetMobile y MobileNetV2), preentrenadas en ImageNet y con la misma cabeza de clasificación, en la distinción binaria
entre melanoma y nevo melanocítico sobre imágenes dermatoscópicas de DermNet. Además de la exactitud y el \gls{auc}
miden el tiempo de inferencia y el tamaño del modelo, y evalúan la significancia de las diferencias con la prueba de
McNemar. DenseNet121 logra la mayor exactitud (92.30\,\%) y ResNet50V2 el mayor \gls{auc} (0.957), pero MobileNetV2
llega a 92.19\,\% con el menor tamaño (9.89 MB) y la inferencia más rápida (23.46 ms), lo que la perfila para
dispositivos de pocos recursos. No hubo expertos ni una referencia externa. Los autores reconocen que la tarea es
binaria, que no cubre nevos atípicos ni otras lesiones pigmentadas, que las imágenes se reducen a 224$\times$224
píxeles y que el modelo no incorpora metadatos del paciente.

\textcite{aboulmira2024skinhealthmate} integran en la aplicación web y móvil SkinHealthMate un ensamble que promedia
las probabilidades de EfficientNetB1 y EfficientNetB5, con preprocesamiento (filtro de mediana) y aumento de datos,
para clasificar las siete categorías de HAM10000. La aplicación permite a los dermatólogos cargar la imagen, validar
el diagnóstico y dar seguimiento al paciente. Reportan una exactitud de 93\,\% y un F1 macro de 0.81, y atribuyen la
diferencia entre ambos a las clases con pocas muestras: en melanoma, la precisión es 0.63 y la sensibilidad 0.49. No
comparan el ensamble con otros métodos ni con clínicos. Reconocen que el modelo, entrenado con HAM10000, no garantiza
un diagnóstico correcto en pieles oscuras, que la calidad de la imagen afecta el resultado y que las condiciones
raras pueden clasificarse mal, por lo que lo presentan como complemento del criterio médico.

\textcite{khan2025ensemble}, publicado solo en arXiv, sin revisión de pares, combinan un ensamble heterogéneo de
EfficientNetB3, ResNet50 y DenseNet121 con un modelo de lenguaje que redacta un reporte a partir de la predicción. El
clasificador distingue nevo de carcinoma basocelular en imágenes dermatoscópicas, con partición estratificada al azar
80/10/10, y decide por mayoría de votos; los casos en que los modelos discrepan se marcan para revisión de un
especialista. Se compara contra cada red por separado con exactitud, F1 y pérdida logarítmica: el ensamble alcanza una
exactitud de 97.67\,\% contra 93.83\,\% de ResNet50, el mejor modelo individual. El texto lo genera LLaMA-3 70B con un
\emph{prompt} que recibe la enfermedad, la confianza y si hubo consenso, y se valora con una rúbrica propia sin
comparar con reportes clínicos. No participaron expertos. Los autores reconocen que el marco es binario, que la
diversidad del conjunto limita el desempeño en distintos tonos de piel y que falta una validación en entorno real.

\textcite{kim2025skinsavvy2} reemplazan la DenseNet de su sistema anterior por un transformer Swin o una EfficientNetV2
y amplían HAM10000 con imágenes de revistas médicas, el archivo ISIC y DermNet NZ para atenuar el desbalance, sin
aumento de datos para no alterar el color ni la textura de la lesión. Las exactitudes por clase del modelo entrenado con
la versión ampliada superan las del entrenado con HAM10000 en casi todas las clases, salvo en la mayoritaria. Con Swin
se obtienen una exactitud balanceada de 0.780, una precisión promedio de 0.812, una sensibilidad promedio de 0.779 y un F1
promedio de 0.793, contra 0.744, 0.727 y 0.734 del sistema anterior en precisión, sensibilidad y F1; la
comparación con cuatro trabajos previos, entre ellos el de Esteva et al., da una exactitud entre 0.705 y 0.776.
Sobre la salida del clasificador un modelo GPT genera recomendaciones de cuidado personalizadas, que valoraron 20
participantes, dos de ellos dermatólogos; esa evaluación mide las recomendaciones y no el diagnóstico. Los autores
reconocen una brecha de dominio entre el conjunto curado y las imágenes tomadas en entornos no controlados.

Entre 2020 y 2023, la línea pasó de la lesión aislada a las condiciones de atención
primaria~\cite{liu2020deep}, corrigió el desbalance dentro de los desafíos ISIC~\cite{gessert2020skin,yao2022single}
y probó un ensamble en un estudio prospectivo~\cite{marchetti2023prospective}. Esos trabajos muestran que un
promedio global oculta la debilidad en los tumores malignos, que el melanoma exige su propia sensibilidad y que el
\gls{auc} se sostiene fuera del desafío, aunque con una especificidad baja y una mejora modesta sobre los
dermatólogos. Los recientes repiten ensambles y comparaciones entre redes, y
reportan exactitudes altas que el propio F1 macro de \textcite{aboulmira2024skinhealthmate} relativiza. Dos de ellos
agregan un texto generado por un modelo de lenguaje sobre la predicción~\cite{khan2025ensemble,kim2025skinsavvy2}.

Quedan cuatro cuestiones abiertas. Casi todos los trabajos usan redes convolucionales, solo
\textcite{kim2025skinsavvy2} incluyen un transformer, y no hay un protocolo común entre trabajos. Solo
\textcite{gessert2020skin}, por lesión, y \textcite{liu2020deep}, por paciente, describen particiones que mantienen
juntos los casos relacionados; \textcite{yao2022single} seleccionan sobre la prueba y \textcite{khan2025ensemble}
parten al azar. La validación fuera de la fuente es temporal o de un solo centro sin lesiones de fototipos V y
VI, los demás evalúan en particiones o pruebas de la misma fuente, y ninguno se desarrolla ni se valida en
población de América Latina. Por último, los mejores sistemas son ensambles grandes, los dos que generan texto no lo
comparan con reportes clínicos de referencia y solo \textcite{khan2025ensemble} marca los casos discordantes, sin
medir esa abstención.

\FloatBarrier
\subsection{Modelos visión--lenguaje en dermatología}
\label{sec:estado_arte_vlm}

Los \glspl{vlm} llegan a la dermatología por las dos vías de la Sección~\ref{sec:marco_vlm}: la alineación
contrastiva, que permite clasificar sin ajuste y recuperar casos, y la unión de un codificador visual con un modelo
de lenguaje que genera texto. Interesa de cada trabajo qué problema ataca, con qué método, contra qué se comparó y si
hubo expertos. Se presentan en orden cronológico dentro de cada estrategia, y las
Tablas~\ref{tab:marco_sota_vlm} y~\ref{tab:marco_sota_vlm_gen} los resumen.

\begin{table}[!t]
    \centering
    \caption{\glspl{vlm} contrastivos y modelos fundacionales en dermatología. $^\dagger$: sin revisión de pares.}
    \label{tab:marco_sota_vlm}
    \footnotesize
    \setlength{\tabcolsep}{3.5pt}
    \begin{tabularx}{\textwidth}{@{}>{\raggedright\arraybackslash}p{1.9cm}
        c
        >{\raggedright\arraybackslash\hspace{0pt}}X >{\raggedright\arraybackslash\hspace{0pt}}X
        >{\raggedright\arraybackslash\hspace{0pt}}p{2.2cm} >{\raggedright\arraybackslash\hspace{0pt}}X
        >{\raggedright\arraybackslash\hspace{0pt}}p{1.9cm}@{}}
        \toprule
        \textbf{Trabajo} & \textbf{Año} & \textbf{Arquitectura} & \textbf{Datos} & \textbf{Métricas} & \textbf{Resultado principal} & \textbf{Expertos} \\
        \midrule
        \textcite{kim2024transparent} & 2024 & CLIP ajustado de forma contrastiva
            & 105\,550 imágenes con texto de PubMed y libros
            & \gls{auc} por concepto
            & Tono de piel: 0.805 contra 0.736 (CLIP)
            & 2 dermatólogos verifican la recuperación \\
        \filasep
        \textcite{yan2025derm1miccv} & 2025 & DermLIP, tipo CLIP, cuatro codificadores
            & Derm1M, ontología con base de 4 dermatólogos
            & Exactitud, \gls{auc}, \emph{Recall@K}
            & \emph{Zero-shot} 0.561 contra 0.441 (BiomedCLIP)
            & Solo la ontología \\
        \filasep
        \textcite{yan2025makemiccai} & 2025 & CLIP con subtextos, alineación local y ponderación
            & 403\,563 pares imagen--texto
            & Exactitud, \gls{auc}, \emph{Recall@K}
            & \emph{Zero-shot} 0.493 contra 0.442 (CLIP, mismo corpus)
            & Ninguno \\
        \filasep
        \textcite{yan2025vision} & 2025 & ViT-Large auto-supervisado con maestro CLIP; no generativo
            & Más de 2 millones de imágenes, 11 instituciones
            & \gls{auc}, F1 ponderado
            & Mejora en 9 de 10 conjuntos; HAM10000 por 4.7\,\%
            & 3 estudios de lectura \\
        \filasep
        \textcite{yan2026dermfmzero}$^\dagger$ & 2026 & ViT-Large auto-supervisado y contrastivo; no generativo
            & Más de 3 millones de imágenes; 1\,014\,085 pares imagen--texto
            & F1 ponderado, \gls{auc}, sensibilidad macro, \emph{Recall@K}
            & \emph{Zero-shot} 0.733 contra 0.561 (DermLIP); atención primaria con ayuda, de 0.266 a 0.482
            & Tres estudios, con 30, 1\,090 y 34 clínicos \\
        \bottomrule
    \end{tabularx}
\end{table}

\begin{table}[!t]
    \centering
    \caption{\glspl{vlm} generativos en dermatología y su evaluación crítica. $^\dagger$: sin revisión de pares. ID: dentro de la distribución; OOD: fuera de ella.}
    \label{tab:marco_sota_vlm_gen}
    \footnotesize
    \setlength{\tabcolsep}{3.5pt}
    \begin{tabularx}{\textwidth}{@{}>{\raggedright\arraybackslash}p{1.9cm}
        c
        >{\raggedright\arraybackslash\hspace{0pt}}X >{\raggedright\arraybackslash\hspace{0pt}}X
        >{\raggedright\arraybackslash\hspace{0pt}}p{2.2cm} >{\raggedright\arraybackslash\hspace{0pt}}X
        >{\raggedright\arraybackslash\hspace{0pt}}p{1.9cm}@{}}
        \toprule
        \textbf{Trabajo} & \textbf{Año} & \textbf{Arquitectura} & \textbf{Datos} & \textbf{Métricas} & \textbf{Resultado principal} & \textbf{Expertos} \\
        \midrule
        \textcite{li2023llavamed} & 2023 & LLaVA de 7B ajustado en dos etapas
            & 600\,000 pares de PubMed Central; instrucciones de GPT-4
            & Puntaje relativo con GPT-4, exactitud en VQA
            & Supera al supervisado en VQA-RAD y PathVQA cerradas; sin evaluar en piel
            & Ninguno \\
        \filasep
        \textcite{zeng2025mmskin} & 2025 & LLaVA-Med de 7B con LoRA
            & 171\,600 imágenes de libros y nueve conjuntos públicos
            & BLEU-4, BERTScore, exactitud
            & \emph{Zero-shot} Patch16 82.34\,\% contra 68.55\,\% (LLaVA-Med)
            & Ninguno \\
        \filasep
        \textcite{shen2025skingptr1}$^\dagger$ & 2025 & Vision-R1-7B congelado con adaptadores
            & 334\,618 muestras con razonamiento sintético
            & Exactitud, puntajes de 0 a 5, peor grupo
            & 82.50\,\% en Derm12345; media de 3.6 sobre 5
            & 5 dermatólogos, 1\,000 casos \\
        \filasep
        \textcite{ru2026dermogpt}$^\dagger$ & 2026 & Qwen3-VL-8B con LoRA y aprendizaje por refuerzo
            & 772\,675 instrucciones; 33\,999 preguntas
            & Exactitud, juez automático, equidad
            & ID 78.04\,\% y OOD 66.48\,\% (Gemini-2.5-Flash: 63.37\,\% y 61.21\,\%)
            & 2 dermatólogos revisan; referencia humana en 100 casos por tarea \\
        \filasep
        \textcite{jiang2026multimodal}$^\dagger$ & 2026 & Evaluación sin ajuste de cinco modelos
            & Tres conjuntos públicos y 5\,811 consultas de un hospital
            & Exactitud de los tres primeros diagnósticos, triaje
            & Imagen sola en consulta real: 1.50\,\% a 13.35\,\% (abiertos), 24.65\,\% (GPT-4.1)
            & 3 dermatólogos validan el juez \\
        \bottomrule
    \end{tabularx}
\end{table}

\textcite{kim2024transparent} proponen MONET para anotar conceptos clínicos (eritema, úlcera, pigmentación) sobre
imágenes sin etiquetas manuales, con el fin de auditar conjuntos de datos y modelos. Es una versión de CLIP ajustada
de forma contrastiva con 105\,550 imágenes dermatológicas y sus descripciones, extraídas de PubMed y de libros de
texto. Se compara con CLIP y con un ResNet-50 supervisado con conceptos anotados, sobre SkinCon y Derm7pt. En la
anotación del tono de piel alcanza un \gls{auc} de 0.805 contra 0.736 de CLIP y 0.874 del supervisado, que se evalúa
en condiciones más favorables. Dos dermatólogos verifican la recuperación de imágenes por concepto, y un clasificador
por conceptos sobre MONET supera, en imágenes clínicas, a las líneas de base en malignidad (\gls{auc} 0.815) y melanoma (0.881). Los autores
aclaran que no sirve para diagnosticar, que fallará con conceptos ausentes del entrenamiento y que no pudieron medir
el tono de piel en imágenes dermatoscópicas.

\textcite{yan2025derm1miccv} entrenan la familia DermLIP, modelos de tipo CLIP con la pérdida contrastiva original, sobre
Derm1M, un corpus de alrededor de un millón de pares imagen--texto dermatológicos organizados según una ontología
cuya base crearon cuatro dermatólogos y que se amplió con ayuda de un modelo de lenguaje. El trabajo, presentado en ICCV 2025, compara con CLIP, SigLIP,
CoCa, PMC-CLIP, BiomedCLIP y MONET en ocho conjuntos y cuatro tareas. En clasificación \emph{zero-shot}, la
exactitud promedio sobre cuatro conjuntos es 0.561 contra 0.441 de BiomedCLIP, y 0.588 con un codificador visual
dermatológico (PanDerm). En la anotación \emph{zero-shot} de conceptos supera a MONET en 1.1 puntos de \gls{auc} en promedio, con
un codificador visual más chico. Los autores reconocen el ruido que introducen los textos educativos, y no hubo evaluación con
clínicos.

\textcite{yan2025makemiccai} observan que CLIP trunca los textos a 77 \emph{tokens} y descarta detalle clínico. MAKE
divide cada descripción en oraciones, a las que suma un texto de enfermedad y otro de conceptos generados por un
modelo de lenguaje; alinea esos subtextos con regiones de la imagen y pondera cada uno según su relevancia
diagnóstica. Se preentrena con 403\,563 pares y se compara con CLIP, SigLIP y CoCa entrenados con el mismo
corpus y con BiomedCLIP, PMC-CLIP y MONET. La exactitud \emph{zero-shot} promedio en cinco conjuntos sube de 0.442
(CLIP sobre el mismo corpus) a 0.493, el \gls{auc} de conceptos mejora 1.11 puntos sobre la mejor línea de base y en
recuperación llega a 0.444 de \emph{Recall@50}. No hay evaluación con expertos ni discusión de limitaciones.

\textcite{yan2025vision} proponen PanDerm, un modelo fundacional visual y no generativo. Se preentrena sin etiquetas
con más de dos millones de imágenes de 11 instituciones y cuatro modalidades, con modelado latente enmascarado y
alineación con las características de CLIP como maestro, sobre cuatro GPU H100. Se evalúa en 28 tareas, contra
DINOv2, SwAVDerm, CLIP, MONET, BiomedCLIP y BiomedGPT. Mejora en 9 de 10 conjuntos de diagnóstico (por 4.7\,\% en
HAM10000 y 9.0\,\% en PAD-UFES-20 sobre el siguiente modelo), y con solo el 10\,\% de las etiquetas suele
superar a los demás. Tres estudios de lectura con clínicos muestran que, con imágenes secuenciales de melanoma, su exactitud supera en
10.2 puntos a la del revisor promedio entre 12 clínicos, que con su ayuda la exactitud con dermatoscopía sube 11
puntos y que la exactitud de los tres primeros diagnósticos de profesionales no dermatólogos sube 16.5 puntos en 128
afecciones. Los autores reconocen que cubren unas 200 de más de 1\,000
afecciones y que la evaluación de equidad es incompleta.

\textcite{yan2026dermfmzero}, disponible en arXiv y Research Square, sin revisión de pares, del mismo grupo, busca que un modelo
fundacional visión--lenguaje apoye la decisión clínica sin ajuste por tarea y evaluarlo con clínicos. Es un
ViT-Large no generativo, preentrenado de forma auto-supervisada con más de tres millones de imágenes y después de
forma contrastiva con 1\,014\,085 pares imagen--texto. Agrega autoencoders esparsos que nombran conceptos y permiten
suprimir las neuronas asociadas a artefactos como reglas y pelo. En siete conjuntos \emph{zero-shot} el promedio de
\gls{auc} (conjuntos binarios) y exactitud balanceada (multiclase) es 0.733 contra 0.561 de DermLIP, y en HAM10000 la sensibilidad macro es 0.744 contra 0.512. En tres estudios con
clínicos, 30 médicos de atención primaria con su ayuda pasan de 0.266 a 0.482 de exactitud en los tres primeros
diagnósticos. Los autores reconocen que los estudios son retrospectivos y simulados y que el preentrenamiento cubre
más de 400 de las más de 2\,000 entidades dermatológicas.

La segunda estrategia une un codificador visual con un modelo de lenguaje. \textcite{li2023llavamed} presentan
LLaVA-Med, un asistente biomédico de propósito general. Parte de LLaVA de 7\,000 millones de parámetros y lo ajusta
en dos etapas con figuras de PubMed Central: primero con 600\,000 pares figura--texto, para alinear el vocabulario,
y después con 60\,000 conversaciones generadas por GPT-4 a partir de los textos, en menos de 15 horas con ocho A100. Se evalúa en tres
conjuntos de preguntas visuales de radiología y patología y con un puntaje relativo calculado por GPT-4. Tras el ajuste
supera al estado del arte supervisado en las preguntas cerradas de VQA-RAD y PathVQA, y en las abiertas solo en SLAKE.
No se evalúa en dermatología ni con expertos, y sus autores reconocen alucinaciones y poco razonamiento en profundidad.
Med-Flamingo, PeFoMed y LLaVA-Tri también se evalúan en preguntas visuales de radiología y patología, y
Med-Flamingo además en problemas clínicos de varias especialidades, pero ninguno en
piel~\cite{moor2023medflamingo,he2025pefomed,xie2025medtrinity25m}.

\textcite{zeng2025mmskin} adaptan ese modelo a la piel con SkinVL. Sobre LLaVA-Med de 7\,000 millones de parámetros
entrenan solo la proyección visual y adaptadores LoRA, con un corpus de imágenes de libros de texto y nueve conjuntos
públicos de diagnóstico reformulados como preguntas (171\,600 imágenes en la variante completa). Se compara con cinco
modelos generales y dos modelos médicos de propósito general, entre ellos LLaVA-Med y HuatuoGPT, en descripción, preguntas y
clasificación con ajuste y \emph{zero-shot}. En \emph{zero-shot} la variante completa alcanza 82.34\,\% de exactitud en
Patch16 contra 68.55\,\% de LLaVA-Med, y es el único modelo que supera 50\,\% en HAM10000. Los textos provienen de
libros de dermatólogos, pero no hubo evaluación de las respuestas por expertos ni se discuten limitaciones, más allá de
pedir mayor capacidad de razonamiento.

\textcite{shen2025skingptr1}, publicado solo en arXiv y sin revisión de pares, buscan un modelo que razone de forma
explícita, produzca un reporte estructurado y rinda de forma pareja entre tonos de piel. Mantienen congelado
Vision-R1-7B y entrenan solo el 0.056\,\% de los parámetros, con un sesgo sobre los \emph{logits}, un adaptador de
mezcla de expertos que incorpora un tono de piel estimado y la destilación de PanDerm. El entrenamiento supervisado usa
334\,618 muestras con razonamiento sintetizado por modelos de lenguaje. En siete conjuntos externos, contra cinco
sistemas entre ellos GPT-4o mini y LLaVA-Med, alcanza 82.50\,\% de exactitud en Derm12345 y entre 42.40\,\% y
69.30\,\% en otros tres. Cinco dermatólogos califican 1\,000 casos con una media de 3.6 sobre 5, pero solo sobre las
salidas de este modelo. Los autores reconocen bucles de autovalidación con incertidumbre alta y una latencia que impide
ejecutarlo en dispositivos de borde.

\textcite{ru2026dermogpt}, publicado solo en arXiv y sin revisión de pares, proponen un modelo que siga el orden
clínico de morfología, razonamiento y diagnóstico, y anuncian que publicarán los datos, el \emph{benchmark} y los
pesos tras la aceptación. Reúnen DermoInstruct,
con 772\,675 instrucciones, y DermoBench, con 33\,999 preguntas. Ajustan Qwen3-VL-8B con LoRA y después con aprendizaje
por refuerzo, con una recompensa que combina la exactitud, la similitud jerárquica dentro de la ontología diagnóstica y
la coincidencia de atributos morfológicos. Comparan 16 modelos. En diagnóstico alcanza 78.04\,\% dentro de la
distribución y 66.48\,\% fuera de ella, contra 63.37\,\% y 61.21\,\% de Gemini-2.5-Flash, el mejor modelo comercial.
Dos dermatólogos revisan línea por línea un núcleo de 900 casos, y dos clínicos de dermatología resuelven 100 casos
por tarea como referencia humana. Los autores reconocen que la
brecha con la referencia humana persiste.

La tercera línea mide qué tan lejos están estos modelos de la consulta real. \textcite{jiang2026multimodal}, también
publicado solo en arXiv y sin revisión de pares, y de un grupo independiente de los anteriores, evalúan sin ajuste
cinco modelos (LLaVA-Med v1.5, SkinGPT-4, InternVL-Chat, MedGemma-4B y GPT-4.1) en diagnóstico diferencial y triaje por
gravedad. Usan tres conjuntos públicos y una cohorte de 5\,811 consultas de un hospital, con 46\,405 imágenes, que no
se distribuye. La exactitud de los tres primeros diagnósticos la calcula un juez automático validado contra tres
dermatólogos. En las consultas con la imagen sola los modelos abiertos obtienen entre 1.50\,\% y 13.35\,\% y GPT-4.1
24.65\,\%. Sin contexto, la sensibilidad de triaje queda por debajo de 0.50 en todos, y el texto del clínico remitente
mejora los resultados pero los vuelve sensibles a errores de ese texto. Los autores concluyen que no corresponde
usarlos sin supervisión y reconocen que los casos provienen de un servicio con internación y que son solo cinco
modelos.

La línea contrastiva avanzó de adaptar CLIP con texto de la literatura (MONET) a corpus dermatológicos con ontología
(DermLIP), textos descompuestos por aspecto (MAKE) y preentrenamiento visual a gran escala (PanDerm, DermFM-Zero), y
solo estos dos últimos se evalúan en estudios de lectura con clínicos. La línea generativa pasó de asistentes biomédicos que no se evaluaron en
piel (LLaVA-Med) a modelos adaptados a ella (SkinVL) y a modelos que razonan y producen un reporte estructurado
(SkinGPT-R1, DermoGPT). La evaluación crítica muestra que los \glspl{vlm} médicos de propósito general rinden poco en
piel y que la brecha crece en consultas reales. Quedan cuestiones abiertas.
Ningún trabajo combina un clasificador supervisado con un generador ni recupera casos parecidos para anclar el texto.
La incertidumbre aparece como limitación, o en DermoGPT para ponderar sus predicciones, y no como una decisión
entre afirmar, ofrecer diferenciales o derivar (Sección~\ref{sec:marco_tres_salidas}). El cómputo es alto, de ocho RTX 4090 para ajustar SkinGPT-R1 a cuatro H100 para
preentrenar PanDerm o DermFM-Zero. Y los estudios con clínicos son pocos: solo PanDerm y DermFM-Zero miden el efecto de
la ayuda sobre sus decisiones, y en los generativos los expertos son dos o tres, o evalúan solo las salidas del propio modelo.

\FloatBarrier
\subsection{Generación de reportes dermatológicos}
\label{sec:estado_arte_reportes}

Un reporte dermatológico reúne la descripción de lo que se ve, un diagnóstico y, cuando corresponde, una
recomendación. Los trabajos difieren en cómo lo producen (un único \gls{vlm} o un clasificador que precede a un
redactor), en la evidencia que usan (solo la imagen o también un contexto recuperado) y en cómo evalúan el texto
(Sección~\ref{sec:marco_rep_ejes}). La Tabla~\ref{tab:marco_generadores} resume los generadores, y la
Tabla~\ref{tab:marco_sota_reportes}, los trabajos que se ocupan de la evaluación y de la recuperación.

\begin{table}[!t]
    \centering
    \caption{Generadores de texto dermatológico. «Métricas con referencia»: BLEU, ROUGE y BERTScore; «---»: no
    informado.}
    \label{tab:marco_generadores}
    \small
    \begin{tabularx}{\textwidth}{@{}>{\raggedright\arraybackslash}p{3.2cm}
        >{\raggedright\arraybackslash\hspace{0pt}}X >{\raggedright\arraybackslash\hspace{0pt}}X l@{}}
        \toprule
        \textbf{Trabajo} & \textbf{Modelo que redacta} & \textbf{Evaluación del texto} & \textbf{Hardware} \\
        \midrule
        SkinGPT-4~\cite{zhou2024pretrained} & MiniGPT-4 con Llama-2-13B & dermatólogos sobre 150 casos reales & 8 A100 \\
        \filasep
        SkinVL~\cite{zeng2025mmskin} & LLaVA-Med con LoRA & métricas con referencia & 8 RTX 3090 \\
        \filasep
        DermatoLlama~\cite{yilmaz2025resourceefficient} & Llama 3.2 11B con LoRA & exactitud y métricas con referencia & 4 RTX 4090 \\
        \filasep
        Skin-R1~\cite{liu2025skinr1} & Qwen2.5-VL-7B & exactitud del diagnóstico en opción múltiple & 2 A100 \\
        \filasep
        SkinGPT-R1~\cite{shen2025skingptr1} & Vision-R1-7B congelado con adaptadores & puntajes de un juez automático y de 5 dermatólogos & 8 RTX 4090 \\
        \filasep
        Khan et al.~\cite{khan2025ensemble} & Llama-3 70B por API & --- & --- \\
        \filasep
        SkinSavvy2~\cite{kim2025skinsavvy2} & GPT-4 & encuesta a 20 participantes, 2 dermatólogos & 4 RTX A6000 \\
        \filasep
        Chakrabarti et al.~\cite{chakrabarti2025skin} & LLaVA-1.5-13B destilado de GPT-4V & métricas con referencia, corrección diagnóstica y revisión de un médico & --- \\
        \filasep
        MpoxVLM~\cite{cao2024mpoxvlm} & Llama-2-7B con LoRA & exactitud del diagnóstico & 2 RTX 4090 \\
        \bottomrule
    \end{tabularx}
\end{table}

\begin{table}[!t]
    \centering
    \caption{Evaluación de narrativas dermatológicas y recuperación de contexto para \glspl{vlm} médicos.}
    \label{tab:marco_sota_reportes}
    \footnotesize
    \setlength{\tabcolsep}{3.5pt}
    \begin{tabularx}{\textwidth}{@{}>{\raggedright\arraybackslash}p{1.9cm}
        c
        >{\raggedright\arraybackslash\hspace{0pt}}X >{\raggedright\arraybackslash\hspace{0pt}}X
        >{\raggedright\arraybackslash\hspace{0pt}}p{2.2cm} >{\raggedright\arraybackslash\hspace{0pt}}X
        >{\raggedright\arraybackslash\hspace{0pt}}p{1.9cm}@{}}
        \toprule
        \textbf{Trabajo} & \textbf{Año} & \textbf{Arquitectura} & \textbf{Datos} & \textbf{Métricas} & \textbf{Resultado principal} & \textbf{Expertos} \\
        \midrule
        \textcite{shen2025dermbench} & 2025 & Juez automático y evaluador LLaVA v1.5 7B con LoRA
            & 4\,000 imágenes de DermNet con referencia; 500 imágenes y nueve modelos
            & Error absoluto medio contra médicos, seis dimensiones
            & Error medio 0.363 (DermBench) y 0.178 (DermEval)
            & 2 dermatólogos, coautores \\
        \filasep
        \textcite{xia2025mmedrag} & 2025 & LLaVA-Med-1.5 con recuperación de reportes y optimización de preferencias
            & Radiología, oftalmología y patología
            & Exactitud, F1, \gls{auc}; BLEU, ROUGE-L, METEOR
            & IU-Xray: exactitud 75.47 a 89.54; BLEU 9.64 a 31.38
            & No \\
        \filasep
        \textcite{xia2024rule} & 2024 & LLaVA-Med-1.5 con recuperación de reportes, $k$ calibrado y optimización
            de preferencias
            & Radiología y oftalmología
            & Exactitud, F1; BLEU, ROUGE-L, METEOR
            & IU-Xray: exactitud 75.47 a 87.84; BLEU 9.64 a 27.53
            & No \\
        \bottomrule
    \end{tabularx}
\end{table}

Un primer grupo genera el texto con un único \gls{vlm} ajustado a dermatología. SkinVL y SkinGPT-R1 pertenecen a
este grupo y se describieron en la Sección~\ref{sec:estado_arte_vlm}; aquí se detallan otros tres.

SkinGPT-4 busca que una persona sin acceso a un dermatólogo suba una foto y obtenga en una conversación la
descripción de la lesión, un diagnóstico y recomendaciones. Ajusta MiniGPT-4 con Llama-2-13B en dos pasos: primero
alinea la imagen con conceptos clínicos y después aprende diagnósticos a partir de notas de médicos, con unas 50\,000
imágenes públicas y propias~\cite{zhou2024pretrained}. Lo compara con MiniGPT-4 sin adaptar y con
cada paso por separado, en ejemplos cualitativos. El texto lo evalúan cinco dermatólogos certificados con un cuestionario
sobre 150 casos reales, sin métricas con referencia: consideran correcto o pertinente el 80.63\,\% de los
diagnósticos. El entrenamiento llevó unas 24 horas en 8 A100 de 80\,GB. Los autores reconocen que los tonos
Fitzpatrick V y VI están poco representados, que las enfermedades con varias lesiones combinadas siguen siendo
difíciles y que las alucinaciones de los modelos de lenguaje pueden dar información errónea al paciente.

DermatoLlama apunta a un modelo que se ejecute localmente, por privacidad y por costo, y lo logra con datos
sintéticos~\cite{yilmaz2025resourceefficient}. A partir de 367 tareas semilla revisadas por expertos y 82\,379
imágenes abiertas, modelos más grandes consultados por API generan 1.2 millones de textos sintéticos (DermaSynth),
con los que ajusta con LoRA un Llama~3.2 de 11 mil millones de parámetros. Se compara con el modelo base, GPT-4o y
Gemini~2.0~Flash en tres evaluaciones. La primera, de 225 imágenes de DERM12345, la anotaron tres dermatólogos
coautores con preguntas jerárquicas; allí DermatoLlama alcanza una exactitud de 0.87 (GPT-4o, 0.86), un BLEU-4 de
0.66 (0.55) y un ROUGE-L de 0.93 (0.81). Se mide, pues, el diagnóstico y el parecido con un texto de referencia, y
ningún experto revisa los textos generados. Entrena en 4 RTX 4090 y ejecuta en una. Los autores advierten que los
datos sintéticos pueden perder rasgos sutiles en casos raros y que faltan la validación clínica externa y la
auditoría de sesgos. Está publicado solo en medRxiv, sin revisión de pares.

Skin-R1 se ocupa de que el diagnóstico venga con un razonamiento fundado, no de redactar un reporte
completo~\cite{liu2025skinr1}. Sobre Qwen2.5-VL-7B con LoRA, genera trayectorias de razonamiento a partir de libros
de texto, con la jerarquía de enfermedades y el diagnóstico diferencial; las usa para el ajuste supervisado y luego
para un aprendizaje por refuerzo con recompensas por exactitud jerárquica, por discriminación de malignidad y
por formato. Lo evalúa con preguntas de
opción múltiple sobre seis conjuntos de datos de dermatología y un conjunto externo de dermatoscopía. Se compara con
\glspl{vlm} generales (Qwen, InternVL, LLaVA), médicos (LLaVA-Med, HuatuoGPT-Vision, MedGemma) y SkinVL, que casi
nunca produce respuestas válidas con ese formato. La exactitud media es 0.6385 en las pruebas de los mismos seis conjuntos
del entrenamiento (0.1955 sobre el siguiente modelo) y 0.7171 en el conjunto externo; en la clasificación benigno, maligno o precanceroso, 0.6928 de
exactitud y 0.4287 de F1 macro. Entrena en 2 A100 de 40\,GB. No mide la calidad del texto del razonamiento ni
incluye expertos, y está publicado solo en arXiv, sin revisión de pares.

Un segundo grupo separa el diagnóstico de la redacción: un clasificador decide y un modelo de lenguaje escribe.
Los dos sistemas de este grupo descriptos en la Sección~\ref{sec:estado_arte_diagnostico} redactan con un
modelo grande. En el de \textcite{khan2025ensemble}, Llama-3 de 70 mil millones de parámetros, consultado por API,
recibe solo el nombre de la enfermedad, el tipo de consenso entre los clasificadores y la confianza, no la imagen, y
redacta la descripción de la lesión, los síntomas por vigilar, el tratamiento y las señales de alarma; el texto solo
se valora con una rúbrica propia de contenido, sin métricas con referencia ni expertos. En SkinSavvy2, la clase predicha se combina en un \emph{prompt} con la edad,
el sexo, el tipo de piel, la exposición solar y los antecedentes familiares, y GPT-4 redacta una recomendación de
cuidado~\cite{kim2025skinsavvy2}. Veinte participantes, 12 legos, 6 profesionales de la salud y 2 dermatólogos,
puntuaron de 0 a 5 la confianza, la personalización, la utilidad y la guía de manejo de 20 respuestas, y el
\emph{prompt} con datos personales y GPT-4o obtuvo los puntajes más altos. Los dermatólogos lo consideraron una ayuda
y no un reemplazo del diagnóstico profesional, y los autores reconocen que los datos que ingresa el usuario pueden ser
inexactos.

Chakrabarti et al.\ generan reportes estructurados (hallazgos, impresión y diagnóstico diferencial) de imágenes
dermatoscópicas según la lista de 7 puntos para melanoma~\cite{chakrabarti2025skin}. Sin pares imagen-reporte
disponibles, GPT-4V actúa como maestro: a partir de la imagen, una descripción de libro de texto, los rasgos anotados
de la lista y el diagnóstico, produce pseudorreportes con los que se ajusta LLaVA-1.5-13B, sobre 305 imágenes. Se
compara con Claude~3.7 Sonnet, con LLaVA sin ajustar y con dos variantes de recuperación (BioMedCLIP y PubMedCLIP)
que anteponen al \emph{prompt} los pseudorreportes más parecidos. Contra los pseudorreportes de referencia, el modelo
obtiene SacreBLEU 55.59, ROUGE-1 0.5438 y BERTScore F1 0.9025; la recuperación con BioMedCLIP, 21.84, 0.5162 y
0.8966. Evalúa además los rasgos de la lista y la sensibilidad y especificidad para melanoma, y un médico revisa 30
reportes de este modelo y de GPT-4V. Como la referencia la escribe el propio maestro, las métricas miden el parecido
con GPT-4V. Los autores atribuyen el SacreBLEU inusualmente alto a la estructura de plantilla de los
pseudorreportes y señalan que las métricas léxicas no bastan. Reconocen también el desbalance de clases y las confusiones entre lesiones pigmentadas benignas y melanoma.

SkinGEN acompaña el diagnóstico con imágenes de referencia generadas por difusión~\cite{lin2025skingen}, y MpoxVLM
combina la imagen con datos clínicos para diagnosticar una sola enfermedad, la viruela
símica~\cite{cao2024mpoxvlm}.

Dos líneas complementan a los generadores: cómo evaluar el texto y cómo condicionar la generación con contexto
recuperado.

\textcite{shen2025dermbench} abordan cómo evaluar las narrativas diagnósticas más allá de las métricas de similitud
textual, que según los autores correlacionan poco con la utilidad clínica. DermBench tiene una narrativa de
referencia para cada una de 4\,000 imágenes de DermNet, redactada por modelos de lenguaje a partir de la etiqueta
verdadera y corregida por dermatólogos. Un juez automático que no ve la imagen compara el candidato con la
referencia y lo puntúa de 0 a 5 en seis dimensiones. DermEval es un evaluador sin referencia, un LLaVA v1.5 7B con
LoRA entrenado con aprendizaje por refuerzo contra los puntajes de médicos. Sobre nueve modelos generales o médicos,
el error absoluto medio contra los médicos, promediado en las seis dimensiones, es 0.363 para DermBench y 0.178
para DermEval. Dos dermatólogos,
coautores, puntuaron todo, sin cuantificar el acuerdo entre ellos. Los autores reconocen que las referencias se
redactan con modelos, lo que puede introducir sesgo. El trabajo se presentó en ECCV 2026.

MMed-RAG~\cite{xia2025mmedrag} y RULE~\cite{xia2024rule}, del mismo grupo, son los antecedentes de la recuperación
para \glspl{vlm} médicos. Buscan mejorar la factualidad de LLaVA-Med-1.5 con contexto recuperado y evitar que el
modelo ignore la imagen o que un contexto equivocado degrade respuestas correctas. En los dos la consulta es la
imagen y lo recuperado son reportes de texto, que se concatenan en el \emph{prompt}; ninguno recupera imágenes ni se
abstiene. RULE calibra cuántos reportes recuperar con un control estadístico del riesgo, y MMed-RAG corta la lista
donde la similitud cae de forma brusca. Los dos ajustan además el generador con optimización directa de preferencias,
para que use el contexto solo cuando ayuda. Los datos son de radiología, oftalmología y, en MMed-RAG, patología; no
hay dermatología. Se comparan con métodos de decodificación y con otros \glspl{vlm} médicos, y MMed-RAG también
con métodos de recuperación, entre ellos RULE. En IU-Xray la exactitud
de las preguntas sube de 75.47 a 87.84 con RULE y a 89.54 con MMed-RAG. En la ablación de RULE, sin embargo, el
contexto recuperado sin más empeora MIMIC-CXR, de 75.79 a 67.35. Ninguno incluye expertos, y MMed-RAG señala que una
modalidad nueva puede exigir reentrenar el recuperador.

El texto se evalúa de tres maneras que casi no se combinan (Sección~\ref{sec:marco_eval_texto}). Las métricas con
referencia miden el parecido con un texto y no la corrección clínica, y las referencias suelen ser sintéticas:
DermatoLlama usa preguntas anotadas por sus coautores y Chakrabarti et al.\ los pseudorreportes del propio maestro.
La corrección del diagnóstico (Skin-R1, DermatoLlama) deja fuera la descripción y la recomendación. Los expertos
son pocos: cinco dermatólogos sobre 150 casos en SkinGPT-4, dos dermatólogos entre 20 participantes en SkinSavvy2,
un médico sobre 30 reportes en Chakrabarti et al.\ y dos coautores en DermBench; Khan et al.\ solo aplican una
rúbrica propia.
La recuperación puede mejorar la factualidad y también degradarla: en la ablación de RULE el contexto sin más
pierde 8.4 puntos en MIMIC-CXR, y en dermatología la recuperación de pseudorreportes rinde menos que el ajuste. En todos los casos
se recuperan textos y no imágenes, y las métricas miden el parecido con una referencia. El hardware va de dos GPU de
consumo a ocho A100 (DermatoLlama se ejecuta en una sola), y varios sistemas dependen de modelos cerrados o consultados por API
para redactar o para generar sus datos o referencias (Llama-3 70B en Khan et al., GPT-4 en SkinSavvy2, GPT-4V en
Chakrabarti et al., Gemini~2.0~Flash en DermatoLlama y Gemini entre otros modelos en DermBench).

\FloatBarrier
\subsection{Posicionamiento del trabajo}
\label{sec:estado_arte_posicion}

Cada línea deja algo sin resolver. Los clasificadores de lesiones descansan sobre la prueba de un desafío, un solo
centro o una partición temporal, y los de mejor desempeño son ensambles grandes. Entre los \glspl{vlm}, los
contrastivos especializados clasifican y recuperan pero no redactan, los generativos exigen un cómputo elevado y los
\glspl{vlm} médicos de propósito general rinden poco en piel. En los reportes, la mayoría de los sistemas se evalúa por parecido textual o por
la exactitud del diagnóstico, y la recuperación estudiada trae textos de otras especialidades.

De esos huecos sale el \emph{pipeline}: un clasificador que estima el diagnóstico, una recuperación de casos
parecidos y un generador de reportes, con una decisión que depende de la confianza y determina cuánto afirma el
texto. Cada pieza atiende una brecha de la Tabla~\ref{tab:marco_brechas}
(Sección~\ref{sec:marco_tres_salidas}), y el análisis incluye cuánto degrada un vecino equivocado el
texto, como advierten MMed-RAG y RULE~\cite{xia2025mmedrag,xia2024rule}.

\begin{table}[!t]
    \centering
    \caption{De las brechas de la revisión a las decisiones del trabajo.}
    \label{tab:marco_brechas}
    \small
    \begin{tabularx}{\textwidth}{@{}>{\raggedright\arraybackslash}p{3.3cm}
        >{\raggedright\arraybackslash\hspace{0pt}}X >{\raggedright\arraybackslash\hspace{0pt}}X@{}}
        \toprule
        \textbf{Brecha} & \textbf{Evidencia en la revisión} & \textbf{Decisión del trabajo} \\
        \midrule
        Sin comparación bajo un mismo protocolo
            & Clasificadores casi solo convolucionales, validados sobre pruebas de desafíos o en un centro
              (Sección~\ref{sec:estado_arte_diagnostico}); \glspl{vlm} medidos entre sí o con clínicos
              (Sección~\ref{sec:estado_arte_vlm})
            & Comparar clasificadores supervisados y modelos médicos generales y dermatológicos con un mismo
              protocolo, con partición disjunta por lesión y un conjunto externo \\
        \filasep
        Texto evaluado por su parecido con una referencia
            & Métricas con referencia, exactitud del diagnóstico o un juez que no ve la imagen
              (Sección~\ref{sec:estado_arte_reportes})
            & Medir por separado la fidelidad diagnóstica y la calidad descriptiva
              (Sección~\ref{sec:marco_rep_ejes}) \\
        \filasep
        Incertidumbre y evidencia sin integrar
            & La incertidumbre como limitación y no como decisión; contexto recuperado de otras especialidades
              (Secciones~\ref{sec:estado_arte_vlm} y~\ref{sec:estado_arte_reportes})
            & Recuperar casos parecidos por imagen, usar la concordancia como señal de confianza y decidir entre
              afirmar, ofrecer diferenciales o derivar
              (Secciones~\ref{sec:marco_integracion} y~\ref{sec:marco_tres_salidas}) \\
        \filasep
        Cómputo fuera del alcance de un equipo sin infraestructura
            & Ocho RTX 4090, cuatro H100, ensambles de 90 redes o un modelo de 70 mil millones de parámetros
              por API
            & Ejecutar todo en una GPU de consumo, con un codificador contrastivo dermatológico y LoRA \\
        \bottomrule
    \end{tabularx}
\end{table}

El codificador contrastivo es DermLIP. DermFM-Zero lo supera en su propio protocolo, pero DermLIP, en su variante
ViT-B/16, se ejecuta en una GPU de consumo, y un mismo espacio sirve para clasificar sin ajuste, entrenar un
clasificador lineal y recuperar vecinos visuales (Sección~\ref{sec:marco_recuperacion_imagenes},
\cite{yan2025derm1miccv}). Los modelos generativos se adaptan con LoRA (Sección~\ref{sec:marco_llm_lora}).

Dos trabajos se acercan a una parte del diseño. El ensamble de \textcite{khan2025ensemble} deriva a un especialista
los casos en que sus miembros no coinciden, pero la concordancia solo agrega al texto un aviso de revisión, y el
texto lo escribe un modelo consultado por API. \textcite{chakrabarti2025skin} miden la corrección diagnóstica de los reportes de un único
modelo destilado. Ninguno combina en una cascada la recuperación de casos, la decisión por concordancia y la
redacción. La implementación y el protocolo de evaluación se describen en el Capítulo~\ref{cp:metodos}.

\FloatBarrier
